% ===========================================================================
\chapter{Desenvolvimento do Trabalho}\label{cap:desenvolvimento}
% ===========================================================================

Este capítulo apresenta a transformação dos requisitos do
\autoref{cap:especificacao} no produto: as tecnologias utilizadas, o projeto e
a implementação da solução, com as decisões tomadas e as alternativas
descartadas, e os testes e a avaliação. As fases 4 e 5 do método ainda estão
em andamento, e o capítulo registra o estado em que se encontram.

\section{Tecnologias Utilizadas}\label{sec:tecnologias}

Esta seção compara, para cada decisão de projeto, a alternativa escolhida com
as descartadas, usando os mesmos critérios para todas. A comparação é
qualitativa: registra propriedades verificadas no código-fonte ou na
documentação de cada alternativa e marca como ``não verificado'' o que não foi
confirmado.

\subsection{Biblioteca de OT}

Critérios:
(C1) compatibilidade com o fluxo de compra, em que o programa local envia a
requisição e o sistema web responde em uma única mensagem;
(C2) possibilidade de guardar ou reproduzir o estado do receptor entre a
requisição e a extração das chaves (RF11, RF14);
(C3) execução no Windows, ambiente do programa local;
(C4) maturidade e recomendação de uso pelos autores;
(C5) aderência ao foco do trabalho, que é o uso de protocolos de OT.

\begin{table}[H]
\centering
\caption{Comparativo das bibliotecas de OT avaliadas}\label{tab:comp-libs}
\footnotesize
\renewcommand{\arraystretch}{1.25}
\begin{tabularx}{\textwidth}{@{}>{\raggedright\arraybackslash}p{0.13\textwidth}
  >{\raggedright\arraybackslash}X >{\raggedright\arraybackslash}X
  >{\raggedright\arraybackslash}p{0.11\textwidth} >{\raggedright\arraybackslash}X
  >{\raggedright\arraybackslash}p{0.11\textwidth}@{}}
\toprule
\textbf{Alternativa} & \textbf{C1 -- Mensagens} & \textbf{C2 -- Estado do receptor} &
\textbf{C3 -- Windows} & \textbf{C4 -- Maturidade} & \textbf{Resultado} \\
\midrule
libOTe, OT base de \citeonline{masnyrindal2019} &
Compatível: a requisição do receptor não depende da mensagem do \emph{sender},
que segue na resposta &
Em memória, mas reproduzível: toda a aleatoriedade do receptor vem do gerador
recebido como parâmetro (verificado no código) &
Sim, segundo os autores &
Acadêmica, licença MIT, amplamente citada &
\textbf{Escolhida} \\
\addlinespace
emp-ot, protocolo PVW &
Compatível: receptor envia primeiro, duas mensagens &
Em memória, chamada bloqueante; reprodutibilidade não verificada &
\textbf{Não} (apenas Linux e macOS) &
Software de pesquisa, sem auditoria independente, segundo os autores &
Descartada (C3) \\
\addlinespace
circl, Simplest OT &
Exige uma troca adicional: a requisição do receptor depende da mensagem inicial
do \emph{sender} &
API dividida em rodadas, favorável &
Sim &
Apenas 1-out-of-2; pacote fora da versão mais recente do módulo &
Descartada (C1, C4) \\
\addlinespace
mpz, Chou--Orlandi &
Mesma dependência da mensagem inicial do \emph{sender} (protocolo Chou--Orlandi) &
Favorável: máquinas de estado sem rede, com mensagens explícitas &
Não verificado &
Autores recomendam não usar em produção; mudanças incompatíveis frequentes &
Descartada (C4) \\
\addlinespace
swanky (Naor--Pinkas, Chou--Orlandi) &
Não verificado &
Não verificado &
Não verificado &
Autores recomendam não usar em produção &
Descartada (C4) \\
\addlinespace
OPRF (RFC 9497) usado para construir 1-out-of-n &
Compatível: cliente envia primeiro, duas mensagens &
API em passos, favorável &
Sim &
Padronizado, várias implementações independentes &
Descartada (C5) \\
\bottomrule
\end{tabularx}
\fonte{Elaborado pelos autores a partir da documentação e do código-fonte das bibliotecas.}
\end{table}

A alternativa baseada em função pseudoaleatória oblívia (\emph{Oblivious Pseudorandom Function} -- OPRF), padronizada no documento RFC (\emph{Request for Comments}) 9497, é tecnicamente adequada, mas não é um protocolo
de OT: ela obtém a funcionalidade de 1-out-of-n a partir de uma função
pseudoaleatória avaliada obliviamente, o que deslocaria o foco do trabalho.


\subsection{Tecnologias do sistema web}

\begin{table}[H]
\centering
\caption{Comparativo das tecnologias do sistema web}\label{tab:comp-stack}
\footnotesize
\renewcommand{\arraystretch}{1.25}
\begin{tabularx}{\textwidth}{@{}>{\raggedright\arraybackslash}p{0.14\textwidth}
  >{\raggedright\arraybackslash}p{0.2\textwidth} >{\raggedright\arraybackslash}p{0.2\textwidth}
  >{\raggedright\arraybackslash}X@{}}
\toprule
\textbf{Tema} & \textbf{Escolhida} & \textbf{Descartada} & \textbf{Motivo} \\
\midrule
Linguagem & Go & TypeScript (Node.js); Python &
Go oferece assinatura cega RSA padronizada \cite{rfc9474} em biblioteca
mantida, biblioteca padrão com servidor HTTP (\emph{Hypertext Transfer Protocol}) e criptografia, e binário único.
Em Python não foi identificada implementação consolidada da RFC 9474. \\
\addlinespace
Banco de dados & PostgreSQL & SQLite &
Escritas concorrentes representativas para o teste de carga (RNF05). \\
\addlinespace
Interface & HTML (\emph{HyperText Markup Language}) gerado no servidor & Aplicação de página única (SPA) &
Menor superfície de ataque e nenhuma dependência de código de terceiros no
navegador. \\
\addlinespace
Busca (RF04) & Pacote completo, navegado no programa local & Vitrine no
navegador com filtro local; busca no servidor sem registros &
Na vitrine no navegador, o filtro é código fornecido pelo sistema web, e a
proteção dependeria de ele ser honesto; na busca no servidor, o sistema web
observa as consultas e a proteção dependeria apenas de disciplina
operacional. \\
\addlinespace
Acesso à biblioteca de OT & Binário auxiliar & Ligação direta ao código C++ &
Isola o código C++ e desacopla a compilação da biblioteca da compilação do
sistema web. \\
\bottomrule
\end{tabularx}
\fonte{Elaborado pelos autores.}
\end{table}


\section{Projeto e Implementação}\label{sec:projeto}

\subsection{Visão geral da solução}

A partir dos requisitos, a ideia inicial do projeto é uma plataforma de venda
de chaves de livros criptografados, que são distribuídos por uma CDN. A
solução é composta por três componentes:

\begin{enumerate}
\item \textbf{Sistema web (site):} administração e publicação do catálogo,
preços discretos, pagamento via gateway, execução do OT e assinatura cega de
tokens. A compra não exige conta; apenas o e-mail é solicitado.

\item \textbf{CDN:} armazena os livros cifrados e entrega um livro mediante um
token com assinatura cega válido. Posteriormente, reporta ao site os tokens
utilizados para ser remunerada.

\item \textbf{Programa local:} instalado no computador do cliente, baixa o
catálogo completo e permite buscar e navegar localmente, gera a requisição de
OT e guarda seu segredo, comunica-se com o sistema web e com a CDN, remove o
cegamento dos tokens e descriptografa os livros localmente.
\end{enumerate}

No fluxo de compra, o cliente escolhe um grupo de livros de mesmo preço e
obtém, por meio do OT, as chaves de N itens desse grupo, sem que o site saiba
quais foram escolhidos. O cliente recebe também N tokens com assinatura cega,
com os quais baixa os livros cifrados na CDN, e em seguida os descriptografa
localmente. Como cada download consome um token e a CDN reporta os tokens
utilizados, torna-se possível a remuneração justa da CDN sem revelar ao site
quais livros foram adquiridos.

A execução do OT é síncrona: durante a compra, o programa local envia a
requisição de OT e recebe a resposta em uma única troca de mensagens com o
sistema web, dentro do limite de RNF03. Toda a aleatoriedade do lado receptor
é derivada de uma semente guardada pelo programa local (RF11); por isso, se a
compra for interrompida após o pagamento, o programa reproduz o estado do
receptor a partir da semente e extrai as chaves usando a resposta do OT
recebida por e-mail (RF14). Esse desenho exige um protocolo de OT em que o
sistema web responda em uma única mensagem e em que o estado do receptor seja
reproduzível a partir da semente.

Os livros de um grupo de mesmo preço são organizados em versões imutáveis: a
inclusão, remoção ou alteração de preço de um livro gera uma nova versão do
grupo, e cada compra referencia a versão sobre a qual os índices do OT foram
escolhidos.


\subsection{Decisões de projeto}

\begin{itemize}
\item \textbf{Execução do OT:} síncrona na compra, com o estado do receptor
reproduzível a partir de uma semente guardada pelo programa local (RF11,
RF14).

\item \textbf{Biblioteca de OT:} libOTe \cite{libote}, com o protocolo de
\citeonline{masnyrindal2019} como OT base 1-out-of-2. A compilação no ambiente
do cliente e a reprodutibilidade do receptor a partir da semente são
verificadas por testes automatizados durante a implementação.

\item \textbf{Sistema web:} linguagem Go, banco PostgreSQL, páginas de
administração renderizadas no servidor, sem recursos de terceiros; catálogo
publicado como pacote completo, sem vitrine pública; biblioteca de OT acessada
por um binário auxiliar. A única página pública é a de apresentação e
download, sem cookies, scripts ou recursos de terceiros, que exibe a versão,
o resumo SHA-256 (\emph{Secure Hash Algorithm}) e a chave pública do catálogo atual para que qualquer pessoa
possa conferi-los.

\item \textbf{Navegação e busca:} no programa local, sobre o pacote completo
do catálogo. As capas seguem no pacote como miniaturas de 200 por 300 pixels em
JPEG (\emph{Joint Photographic Experts Group}), geradas no envio da capa, para que o tamanho do pacote não dependa das
imagens originais.

\item \textbf{Pacote do catálogo:} um arquivo ZIP único por versão, com os
metadados em JSON (\emph{JavaScript Object Notation}) e as miniaturas. O arquivo é gerado de forma
determinística, de modo que o mesmo conteúdo produz sempre os mesmos bytes e o
resumo SHA-256 identifica a versão. O programa local sempre baixa o pacote
completo, revalidando-o pelo resumo, para que todos os clientes façam a mesma
requisição. O sistema web assina cada versão com Ed25519 e publica versão,
resumo e assinatura, o que permite a qualquer pessoa verificar que recebeu o
mesmo catálogo que os demais. O programa local confere a assinatura com a
chave pública embutida e recusa versões anteriores à que já possui.

\item \textbf{Assinatura cega:} RSA cego padronizado \cite{rfc9474}, com
módulo de 3072 bits (RNF01) e uma chave por faixa de preço. A chave pública de
cada faixa vai no pacote assinado do catálogo, de modo que todos os clientes e
a CDN recebem a mesma chave e o sistema web não consegue reconhecer os tokens
de um cliente específico.

\item \textbf{Variante de OT:} o OT 1-out-of-n de \citeonline{naorpinkas1999},
repetido N vezes. Cada execução usa $\lceil \log_2 n \rceil$ OTs 1-out-of-2 do
MR19. A chave do livro de índice $i$ é cifrada com um resumo SHA-256 das chaves
dos OTs correspondentes aos bits de $i$. A privacidade da escolha diante do
sistema web é garantia do protocolo; a de que o cliente não abre outras
posições depende de modelar o SHA-256 como oráculo aleatório (premissa).

\item \textbf{Cifra dos livros:} AES-256-GCM, com uma chave aleatória por livro
e o identificador do livro como dado associado.

\item \textbf{Grupos, reenvio e pagamento:} grupo mínimo de 5 livros; reenvio
da resposta do OT por 30 dias; gateway de pagamento simulado atrás de uma
interface, cobrando o preço da faixa vezes N.

\item \textbf{Tokens e reporte:} cada token vale um único download. A CDN
reporta em lotes de 50 tokens ou uma vez por dia, embaralhados, sem horários e
sem livro.

\item \textbf{Semente:} toda a aleatoriedade do cliente numa compra (receptor
do OT, mensagens dos tokens e fatores de cegamento) é derivada de uma semente
de 256 bits gravada antes do envio. Com ela e a resposta recebida por e-mail, o
programa local refaz a compra interrompida.

\item \textbf{Programa local:} em Go, com linha de comando e uma janela gráfica
para Windows que usa o componente WebView2. A página da janela é embutida no
executável e não acessa a rede: todas as operações passam pelo código do
programa, que se conecta apenas ao sistema web e à CDN.

\item \textbf{Remuneração da CDN:} tarifa fixa por token entregue; o sistema
web registra a quantidade e o valor de cada lote.

\item \textbf{Falha seletiva no OT:} o programa local troca todos os tokens na
CDN antes de decifrar e não informa ao sistema web se alguma chave falhou, de
modo que sua comunicação é a mesma em qualquer caso.

\item \textbf{Interface:} um único tema visual para os três componentes, com
fontes servidas pelos próprios componentes, sem recursos de terceiros.
\end{itemize}


\subsection{Alternativas de projeto descartadas}\label{sec:alternativas}

As alternativas a seguir são comparadas com o mesmo método da
\autoref{sec:tecnologias}.

\subsubsection{Execução do OT}

Critérios: utilidade do envio da resposta por e-mail (RF14) e do reenvio
(RF15); recuperação de uma compra interrompida após o pagamento (RNF07);
quantidade de código criptográfico próprio.

\begin{table}[H]
\centering
\caption{Comparativo das formas de execução do OT}\label{tab:comp-execucao}
\footnotesize
\renewcommand{\arraystretch}{1.25}
\begin{tabularx}{\textwidth}{@{}>{\raggedright\arraybackslash}p{0.2\textwidth}
  >{\raggedright\arraybackslash}X >{\raggedright\arraybackslash}X
  >{\raggedright\arraybackslash}X >{\raggedright\arraybackslash}p{0.12\textwidth}@{}}
\toprule
\textbf{Alternativa} & \textbf{RF14 e RF15} & \textbf{Compra interrompida} &
\textbf{Código criptográfico próprio} & \textbf{Resultado} \\
\midrule
Síncrona, com semente guardada &
Úteis: a resposta recebida por e-mail permite extrair as chaves &
Recuperável a partir da semente &
Nenhum além de inicializar o gerador com a semente; depende da premissa de
determinismo do receptor &
\textbf{Escolhida} \\
\addlinespace
Síncrona, sem estado guardado &
Perdem a utilidade: sem o segredo, a resposta não serve &
Irrecuperável: o cliente paga e perde a compra &
Nenhum &
Descartada \\
\addlinespace
Adaptador com estado serializado &
Úteis &
Recuperável a partir do estado salvo &
Reorganização do protocolo em duas etapas sobre as primitivas da biblioteca &
Descartada \\
\bottomrule
\end{tabularx}
\fonte{Elaborado pelos autores.}
\end{table}


\subsubsection{Onde executa o receptor do OT}

Critérios: natureza da privacidade da escolha (RF07), usabilidade, viabilidade
técnica com a biblioteca escolhida e armazenamento da semente (RF11).

\begin{table}[H]
\centering
\caption{Comparativo dos locais de execução do receptor}\label{tab:comp-receptor}
\footnotesize
\renewcommand{\arraystretch}{1.25}
\begin{tabularx}{\textwidth}{@{}>{\raggedright\arraybackslash}p{0.19\textwidth}
  >{\raggedright\arraybackslash}X >{\raggedright\arraybackslash}X
  >{\raggedright\arraybackslash}X >{\raggedright\arraybackslash}X
  >{\raggedright\arraybackslash}p{0.11\textwidth}@{}}
\toprule
\textbf{Alternativa} & \textbf{Privacidade da escolha} & \textbf{Usabilidade} &
\textbf{Biblioteca de OT} & \textbf{Semente} & \textbf{Resultado} \\
\midrule
Programa local com acesso restrito ao sistema web e à CDN &
Garantia do protocolo &
Sem transporte manual de arquivos &
Execução nativa &
Disco do cliente &
\textbf{Escolhida} \\
\addlinespace
Programa local sem rede (versão 2 desta especificação) &
Garantia do protocolo &
Cliente transporta arquivos entre navegador e programa &
Execução nativa &
Disco do cliente &
Descartada \\
\addlinespace
Navegador, com código fornecido pelo sistema web &
Premissa: o servidor fornece código honesto &
Nada a instalar &
Exigiria compilação para WebAssembly (viabilidade não verificada) &
Guardada no navegador, sujeita a remoção &
Descartada \\
\addlinespace
Navegador, com código de outra origem (extensão ou página estática auditável) &
Premissa mais fraca: código auditável e independente do servidor &
Instalação de extensão ou uso de página separada &
Mesma limitação da alternativa anterior &
Guardada no navegador &
Descartada \\
\bottomrule
\end{tabularx}
\fonte{Elaborado pelos autores.}
\end{table}


\subsection{Estado da implementação}\label{sec:estado}

Esta seção registra o andamento da Fase 4 em 22 de novembro de 2026. Os
três componentes estão implementados e o fluxo de compra roda de ponta a
ponta. O código está no
repositório do projeto (\url{https://github.com/migs39/tcc}).

\begin{table}[H]
\centering
\caption{Passos da implementação do sistema web}\label{tab:passos}
\footnotesize
\renewcommand{\arraystretch}{1.2}
\begin{tabularx}{\textwidth}{@{}>{\raggedright\arraybackslash}p{0.075\textwidth}
  >{\raggedright\arraybackslash}X >{\raggedright\arraybackslash}p{0.22\textwidth}
  >{\raggedright\arraybackslash}p{0.15\textwidth}@{}}
\toprule
\textbf{Passo} & \textbf{Entrega} & \textbf{Requisitos} & \textbf{Situação} \\
\midrule
0 & Base: banco de dados, migrações e testes isolados & --- & Concluído \\
1 & Cadastro, autenticação e gestão de administradores & RF01, UC01 & Concluído \\
2 & Catálogo, preços discretos, capas e grupos versionados & RF02, RF03, RF07, UC02, UC03 & Concluído \\
3 & Pacote do catálogo e página pública & RF04, UC04 & Concluído \\
4 & Compra sem conta, e-mail e pagamento & RF05, RF06, RF13 & Concluído \\
5 & Envio da resposta do OT por e-mail & RF14, RNF04 & Concluído \\
6 & OT no sistema web, com a biblioteca libOTe & RF07, RF08 & Concluído \\
7 & Emissão de tokens com assinatura cega & RF16 & Concluído \\
8 & Recebimento e validação do reporte da CDN & RF18, UC17 & Concluído \\
9 & Reenvio da resposta do OT & RF15 & Concluído \\
10 & Testes de privacidade e de carga & RNF02--RNF05 & Parcial \\
\bottomrule
\end{tabularx}
\fonte{Elaborado pelos autores.}
\end{table}


\subsubsection{O que os passos concluídos entregam}

\begin{itemize}
\item \textbf{Administradores (RF01):} cadastro, autenticação e gestão de
perfis, inclusive desativação, que nunca deixa o sistema sem administrador
ativo. As senhas são guardadas com argon2id; a sessão usa um token aleatório
de 256 bits, do qual só o resumo fica no banco. Requisições de alteração
vindas de outra origem são recusadas.

\item \textbf{Catálogo e preços (RF02, RF03):} livros com título, autor,
categoria, descrição, ano, ISBN (\emph{International Standard Book Number}), editora e capa; faixas de preço e categorias
mantidas pelo administrador. A restrição de preço ao conjunto pré-definido é
imposta pelo código e também pelo banco de dados, de modo que nenhum livro
fica com preço fora do conjunto mesmo que o código seja contornado. Remover
um livro o desativa, preservando as versões de grupo que o referenciam.

\item \textbf{Grupos versionados (RF07):} cada mudança na composição de um
grupo de mesmo preço gera uma versão imutável; editar dados de um livro sem
mudar seu preço não gera versão. A posição do livro na versão é o índice
usado pelo OT.

\item \textbf{Pacote do catálogo (RF04):} o sistema web publica o catálogo
como um arquivo ZIP determinístico, com os metadados dos livros ativos, as
versões atuais dos grupos e as miniaturas das capas, assinado com Ed25519. A
página pública de apresentação e download mostra a versão, o resumo e a chave
pública do catálogo atual.
\end{itemize}


\subsubsection{Compra, CDN e programa local}

\begin{itemize}
\item \textbf{Compra (RF05--RF08, RF13, RF16):} o programa local envia o
e-mail, a versão do grupo, N, a requisição do OT e os N valores cegados. O
sistema web roda o lado \emph{sender} do OT, cobra o grupo vezes N, assina os
valores às cegas e devolve a resposta do OT com as assinaturas. Uma
requisição inválida é recusada antes da cobrança.
\item \textbf{E-mail e reenvio (RF14, RF15):} a resposta do OT vai por e-mail,
fora da requisição, e pode ser reenviada ao e-mail original por 30 dias. Ela
não contém chave alguma.
\item \textbf{CDN (RF17, RF18):} guarda os livros cifrados, entrega um livro
por token válido da faixa do livro e ainda não usado e reporta os tokens em
lotes. O sistema web confere cada assinatura e recusa tokens repetidos.
\item \textbf{Programa local (RF04, RF10--RF12):} baixa e confere o catálogo,
faz a busca localmente, grava a semente antes do envio, compra, troca os tokens
por livros na CDN, extrai as chaves e decifra no computador do cliente.
\end{itemize}

O binário auxiliar que expõe o libOTe compila no Windows e no Linux, e a
requisição do receptor é idêntica byte a byte quando gerada com a mesma
semente, o que confirma a premissa da execução síncrona com semente.

\subsection{Questões em aberto}

\begin{itemize}
\item Valor da remuneração da CDN por token (o sistema web já registra a
quantidade por lote).
\item Proteção da semente armazenada no programa local, além da permissão do
arquivo.
\item Compromissos das chaves dos livros no pacote do catálogo, para o
programa local conferir as chaves extraídas (defesa contra o ataque de falha
seletiva no OT).
\item Interface gráfica do programa local.
\item Integração com um gateway de pagamento e um serviço de e-mail reais.
\item Momento do download logo após a compra, que a CDN observa; o programa
local permite adiar o download.
\end{itemize}

% ---------------------------------------------------------------------------

\section{Testes e Avaliação}\label{sec:testes}

\subsection{Testes automatizados}

O sistema web tem 29 testes automatizados, executados contra um banco
PostgreSQL real, cada teste em um esquema isolado. Os grupos de testes são:

\begin{itemize}
\item \textbf{Administradores (9 testes):} formato e verificação do resumo de
senha; cadastro e validações; autenticação; desativação, que encerra sessões
e bloqueia o acesso; troca de senha, que encerra as sessões; sessão guardada
apenas como resumo.
\item \textbf{Catálogo (9 testes):} preço fora do conjunto recusado pelo
código e pelo banco (RF03); regras de faixas e categorias; validação de ano e
ISBN; sequência e imutabilidade das versões de grupo; capas e miniaturas.
\item \textbf{Pacote do catálogo (3 testes):} conteúdo e ordem dos grupos;
determinismo; geração de versões; recusa de pacote adulterado, de número de
versão trocado e de chave errada.
\item \textbf{Interface web (8 testes):} acesso administrativo restrito a
administradores autenticados; recusa de requisições de outra origem; fluxo
completo do catálogo; revalidação do pacote; página pública sem cookies,
scripts ou recursos de terceiros.
\end{itemize}

Os demais componentes acrescentam os seguintes grupos de testes:

\begin{itemize}
\item \textbf{Protocolo (6 testes):} o cliente recupera exatamente as chaves
escolhidas e nenhuma outra posição abre (RF12); a mesma semente gera a mesma
requisição; o tamanho da requisição não depende da escolha; tokens válidos,
reproduzíveis a partir da semente e diferentes do que o sistema web viu
(RF16); cifra dos livros.
\item \textbf{Compra (7 testes):} compra de ponta a ponta com o binário do
libOTe; registros e respostas de mesma forma e tamanho para escolhas
diferentes; e-mail sem chave; cobrança pelo grupo vezes N; pagamento recusado
sem registro; validações; reenvio dentro e fora do prazo; reporte da CDN
válido, repetido e adulterado.
\item \textbf{CDN (1 teste):} entrega só com token da faixa do livro, uma única
vez; reporte em lote; tabela de tokens sem livro nem horário.
\end{itemize}

Um teste de ponta a ponta com os três componentes cobriu o cadastro de livros,
o envio dos arquivos à CDN, a compra de dois livros de um grupo de seis, a
decifragem, as recusas (grupo pequeno, grupos diferentes, pagamento recusado),
a recuperação de uma compra pelo e-mail e o reporte da CDN aceito pelo sistema
web. Falta o teste de carga concorrente para RNF05.


\subsection{Medições}

As miniaturas de capa, de 200 por 300 pixels em JPEG, mediram de 1,6~KB (capa
de cor lisa) a 19,8~KB (pior caso, imagem de ruído aleatório). Com mil livros,
o pacote do catálogo fica abaixo de cerca de 20~MB.

Com um grupo de 100 livros e N = 5, uma compra leva cerca de 58~ms no sistema
web, num computador de dois núcleos. O tempo inclui o \emph{sender} do OT, as
cifras de Naor--Pinkas, a cobrança simulada, cinco assinaturas RSA de 3072
bits e a gravação. Isso dá cerca de mil compras por minuto em um núcleo, acima
das 600 de RNF05, e fica muito abaixo do limite de um minuto de RNF03. A
requisição do OT tem 1.256 bytes para 19 OTs e a resposta do \emph{sender},
64 bytes.
